\documentclass[11pt]{article}
\usepackage{amsmath,amssymb,amsthm,geometry,hyperref}
\geometry{margin=1in}
\newcommand{\per}{\operatorname{per}}
\newcommand{\Sym}{\operatorname{Sym}}
\newtheorem{theorem}{Theorem}
\newtheorem{lemma}[theorem]{Lemma}
\title{Exact tensor-square certificates for two further cases\\of Marcus's block-permanent conjecture}
\author{Research proof for independent verification}
\date{6 November 2026}
\begin{document}
\maketitle
\begin{abstract}
We give computer-assisted, exact rational certificates for the complete cases
$(m,k)=(4,2)$ and $(3,3)$ of Marcus's block-permanent conjecture, including
its equality condition. The certificates establish stronger inequalities for
arbitrary symmetric tensors in specified dimensions. They are verified by two
separate implementations, one using only Node's built-in libraries and exact
integer and rational arithmetic. The arbitrary-$m,k$ conjecture remains
unresolved in this work; novelty has not been independently assessed.
\end{abstract}

\begin{theorem}\label{main}
Let $A$ be a complex Hermitian positive semidefinite matrix, partitioned into
$m\times m$ blocks of size $k\times k$, and let $G_{ij}=\per A_{ij}$.
If $(m,k)=(4,2)$ or $(m,k)=(3,3)$, then
\[
\per A\geq\per G.
\]
When $A$ has no zero row, equality holds exactly when $A$ is block diagonal.
In particular this includes the equality clause with positive definite diagonal
blocks in the original conjecture.
\end{theorem}

\section{The stronger tensor statements}
The Fock inner product satisfies
$\langle z^\alpha,z^\beta\rangle_F=\mathbf1_{\alpha=\beta}\alpha!$
and is conjugate-linear in its first entry. For a symmetric order-$k$ tensor
$Q$ on $\mathbb C^r$, write
\[
p_Q(z)=\sum_{a_1,\ldots,a_k=1}^r Q_{a_1\cdots a_k}z_{a_1}\cdots z_{a_k}.
\]
In particular $\langle p_Q,p_R\rangle_F=k!\langle Q,R\rangle$.
For $m$ such tensors define
\[
F(Q_1,\ldots,Q_m)=\left\|\prod_i p_{Q_i}\right\|_F^2
-\per[\langle p_{Q_i},p_{Q_j}\rangle_F]_{i,j=1}^m.
\]
The support space $U_Q$ is the span of the vectors obtained by fixing all but
one index of $Q$. Equivalently it is the smallest subspace whose $k$th
symmetric power contains $Q$.

\begin{lemma}\label{certified}
The inequality $F\geq0$ holds for arbitrary symmetric tensors for both
parameter triples
\[
(m,k,r)=(4,2,6),\qquad (m,k,r)=(3,3,9).
\]
If every tensor is nonzero, equality holds if and only if their support spaces
are pairwise orthogonal. Both statements also hold in smaller dimension by
zero padding.
\end{lemma}

The exact certificates for this lemma are the accompanying files
\texttt{marcus-invariant-exact-4-6.json} and
\texttt{marcus-invariant-exact-3-9-k3.json}. The next two sections specify
their meaning and give a finite exact proof procedure. These files contain
the rational numbers in the identities, rather than numerical solver output.

\section{Definition and interpretation of the certificates}
Fix a cut $c\in\{1,\ldots,m-1\}$ and set
$L=kc-1$, $R=k(m-c)-1$, and $B=L+R=km-2$.
A network consists of one list of $k$ indices for each of the $m$ tensor
factors, and a list of Kronecker delta pairs. Indices $0,\ldots,B-1$ are free;
the first $L$ are holomorphic free indices and the rest are antiholomorphic.
Every free index occurs once, either in a tensor or in a delta. Internal
indices occur twice, joining one tensor on each side of the cut. Each delta
joins one holomorphic and one antiholomorphic free index. There is at least
one internal contraction; the number of deltas is one less than the number
of internal contractions.

For such a network $n$, $T_n(Q)$ is the tensor obtained by taking
$Q_1,\ldots,Q_c,\overline{Q_{c+1}},\ldots,\overline{Q_m}$, placing their
indices as specified, multiplying by the specified deltas, and summing all
internal indices from $1$ to $r$. The free indices label its entries.
The certificate field \texttt{networks} lists these networks in order for
each cut. All entries are integers; indices, tensor labels and matrix entries
in the JSON files are zero-based.

Each block $b$ of the certificate gives an integer matrix $C_b$ with one row
per network for its cut, and a rational symmetric matrix $H_b$. Set
\[
V_{b,a}=\sum_n(C_b)_{na}T_n,
\qquad S_b(Q)=\sum_{a,d}(H_b)_{ad}\langle V_{b,a},V_{b,d}\rangle.
\]
Here $C_b$ is the field \texttt{basis}, and $H_b$ is \texttt{gram}; rational
entries are written as exact integer fractions. The identity certified is
\begin{equation}\label{identity}
 F(Q)=\frac1{m!}\sum_{\pi\in S_m}\sum_b
 S_b(Q_{\pi(1)},\ldots,Q_{\pi(m)}).
\end{equation}
Every $H_b$ is positive definite, so the right side is nonnegative.
The block names record how the bases were discovered but are not assumptions
needed to verify the inequality. Only the cut prefix has mathematical
meaning for the verifier.

\section{Exact verification of the identity and positivity}
An Einstein contraction with one copy each of $Q_i$ and $\overline{Q_i}$ is
specified by an $m\times m$ table $N$ of nonnegative integers with every row
and column sum $k$. Its invariant $I_N$ contracts $N_{ij}$ indices between
$Q_i$ and $\overline{Q_j}$. Symmetry of the tensors makes the choice of indices
irrelevant. Wick expansion gives
\[
\left\|\prod_i p_{Q_i}\right\|_F^2
=\sum_N\frac{(k!)^{2m}}{\prod_{ij}N_{ij}!}\,I_N.
\]
The terms with $N=kP$, $P$ a permutation matrix, are exactly $\per G$.
Thus $F$ has coefficient zero on those tables and the displayed positive
coefficient on every other table.

To expand $\langle T_n,T_{n'}\rangle$, glue the free indices of the two
networks, conjugating the first. Each remaining index component either joins
one holomorphic tensor leg to an antiholomorphic one, or is a closed loop
of deltas. It therefore equals $r^\ell I_N$, where $\ell$ is the number of
closed loops. This computes every coefficient of every $S_b$ using integers
and the rational entries of $H_b$.

Both sides of \eqref{identity} are invariant under simultaneous permutation
of tensor labels. Their real coefficient expansions are also invariant under
$N\mapsto N^{\mathsf T}$. It is therefore sufficient to compare the total
coefficient in each orbit under these operations. Averaging the right side
by $1/m!$ distributes its orbit total uniformly, exactly as on the left.
The following are the exact verification counts:
\[
\begin{array}{c|r|r|r|r|r}
(m,k,r)&\text{Wick tables}&\text{orbits}&\text{Gram blocks}
&\text{cut-1 networks}&\text{spanning rank}\\\hline
(4,2,6)&282&25&15&360&360\\
(3,3,9)&55&16&11&320&320
\end{array}
\]
Every orbit coefficient agrees exactly. A rational $LDL^{\mathsf T}$
decomposition gives strictly positive pivots for every $H_b$, proving
positive definiteness without an eigenvalue tolerance.

The independent verifier \texttt{marcus-invariant-verify.mjs} enumerates
the individual Wick pairings ($8!$ or $9!$), constructs the glued contraction
tables directly from the saved networks, checks every rational coefficient
equation, and performs rational $LDL^{\mathsf T}$. It uses only Node built-ins.
Run
\begin{verbatim}
node marcus-invariant-verify.mjs marcus-invariant-exact-4-6.json
node marcus-invariant-verify.mjs marcus-invariant-exact-3-9-k3.json
\end{verbatim}
The separate Python verifier \texttt{marcus-invariant-certify.py} reaches
the same conclusion using the factorial formula for Wick coefficients.
Floating-point semidefinite optimization was used to discover candidate
matrices, but it is not used in either final positivity or identity check.

\section{Equality in the tensor statements}
Equality in \eqref{identity}, with every $H_b$ positive definite, forces every
$V_{b,a}$ to vanish for every permutation of the input tensors.
For cut one, the columns of the matrices $C_b$, together with their
permutations of free holomorphic and antiholomorphic boundary positions,
span the entire network space. This is checked exactly: the standalone
verifier closes this span under adjacent boundary transpositions and finds
full rank over the prime field of order $1000003$. Full rank modulo that
prime implies full rank over $\mathbb Q$, and hence over $\mathbb C$.
Boundary permutations act as isometries on tensor entries, so the vanishing
of the $V_{b,a}$ implies vanishing of every cut-one raw tensor $T_n$.

In particular take a network with just one internal contraction, between
$Q_i$ and $\overline{Q_j}$, and no deltas. The other nonzero factors occur
as an ordinary tensor product. Its vanishing therefore says
\[
\sum_s(Q_i)_{s\alpha}\overline{(Q_j)_{s\beta}}=0
\quad\hbox{for every }\alpha,\beta.
\]
By permuting block labels this holds for every pair $i\ne j$. It is precisely
$U_{Q_i}\perp U_{Q_j}$. Conversely pairwise orthogonal support spaces give
disjoint Gaussian coordinates and factorization, proving equality. This
finishes the proof of Lemma~\ref{certified}.

\section{Application to three blocks of size three}
Choose Gram vectors $v_{ia}$ for $A$ and let
$p_i=\prod_{a=1}^kL_{v_{ia}}$, where $L_v(z)=\sum_s v_sz_s$.
Then Gaussian expansion gives
\[
\per A=\left\|\prod_i p_i\right\|_F^2,\qquad
G_{ij}=\langle p_i,p_j\rangle_F.
\]
For $m=k=3$ the nine Gram vectors span at most nine dimensions, so
Lemma~\ref{certified} applies directly.

For a nonzero product of linear forms its support space is exactly the span
of its coefficient vectors. Indeed, its derivative in a direction vanishes
identically if and only if every factor has zero derivative in that direction:
under translation the degree in the translation parameter is the number of
factors with nonzero directional derivative. Thus the equality statement in
the lemma says precisely that the three block spans are orthogonal. This is
block diagonality of $A$.

\section{Application to four blocks of size two}
We use the previously proved three-pairs theorem, with its explicit integer
tensor-square identity, in \texttt{marcus-three-pairs-proof.tex}.
It says that $\|p_1p_2p_3\|_F^2\geq\per\operatorname{Gram}(p_1,p_2,p_3)$
for split quadratics in every dimension, with equality for nonzero factors
exactly when the three block spans are pairwise orthogonal.

Let $U$ be the span of the six vectors in the first three blocks, so
$\dim U\leq6$. In coordinates $z\in U$ and $w\in U^\perp$, expand
\[
p_4(z,w)=p_0(z)+\sum_{\alpha\ne0}h_\alpha(z)
                    \frac{w^\alpha}{\sqrt{\alpha!}}.
\]
Put $f=p_1p_2p_3$ and $q=\per\operatorname{Gram}(p_1,p_2,p_3)$.
Only $p_0$ contributes to inner products with the first three forms. Expanding
the final diagonal entry of the Gram permanent and using orthogonality of
the $w$ monomials gives
\begin{align*}
\|fp_4\|_F^2-\per G
={}&F(p_1,p_2,p_3,p_0)
 +\sum_{\alpha\ne0}(\|fh_\alpha\|_F^2-q\|h_\alpha\|_F^2).
\end{align*}
The first term is nonnegative by the $(4,2,6)$ certificate. The elementary
Fock multiplication inequality $\|fh\|_F^2\geq\|f\|_F^2\|h\|_F^2$,
together with the three-pairs theorem, makes every other term nonnegative.
For completeness that multiplication inequality follows from normal ordering:
\[
\|fh\|_F^2=\sum_\beta\frac1{\beta!}
\| (\partial^\beta f)^*(\partial)h\|_F^2;
\]
the terms of maximum derivative degree give the required product of norms.

Assume all eight vectors are nonzero and equality holds. If $p_0\ne0$,
equality in the certificate would force its support, contained in $U$, to be
orthogonal to the supports of the first three forms, which span $U$.
This is impossible, so $p_0=0$. Some $h_\alpha$ is nonzero; its summand then
forces $\|f\|_F^2=q$, and the first three block spans are pairwise orthogonal.

Write the fourth block vectors as $x+x'$ and $y+y'$, with $x,y\in U$ and
$x',y'\in U^\perp$. Since $p_0=L_xL_y=0$, at least one of $x,y$ is zero.
If $x=0$ and $y\ne0$, nonzeroness of $x'$ gives an outside-degree-one
coefficient $h_\alpha$ which is a nonzero multiple of $L_y$.
The linear multiplication identity
\[
\|fL_y\|_F^2-\|f\|_F^2\|y\|^2=\|D_y f\|_F^2
\]
then makes its summand strictly positive, since the factors of $f$ span $U$
and $y\in U\setminus\{0\}$. The other case is identical. Hence $x=y=0$.
All four block spans are orthogonal, proving the required equality condition.
The converse follows by permanent factorization. This completes the proof
of Theorem~\ref{main}.

\section{Scope}
Together with the earlier three-pairs identity, this proves the complete
cases $(3,2)$, $(4,2)$ and $(3,3)$. The certificates do not establish a
uniform inequality for arbitrary numbers of blocks or arbitrary block sizes.
No resolution of that remaining statement is claimed here.
\end{document}
